% !TEX TS-program = xelatex
% !TEX program = xelatex
% --- UNIVERSAL PREAMBLE BLOCK ---
\documentclass[11pt, a4paper]{article}
\usepackage[a4paper, top=2.5cm, bottom=2.5cm, left=2cm, right=2cm]{geometry}
\usepackage{fontspec}

\usepackage[bidi=bidi, provide=*]{babel}
\babelprovide[import, main]{english}

\babelfont{rm}{Noto Sans}

\usepackage{amsmath}
\usepackage{amssymb}
\usepackage{braket}
\usepackage{booktabs}
\usepackage{tabularx}
\usepackage{hyperref}
\usepackage{enumitem}
\usepackage{parskip}

\hypersetup{
    colorlinks=true,
    linkcolor=blue,
    citecolor=blue,
    filecolor=magenta,
    urlcolor=blue,
    pdftitle={Engineering-as-Code for Quantum ASICs: Unified Whitepaper},
}

\begin{document}

\title{\textbf{Engineering-as-Code for Quantum ASICs: Programmable Metasurfaces, Cryogenic Control, and Applications}}
\author{Kenneth Olsen\\ \small\textit{Prepared with LLM-assisted synthesis.}}
\date{}

\maketitle

\begin{abstract}
The trajectory of quantum information science is constrained by the physical wiring crisis at millikelvin temperatures and by environmental limits on distributing entanglement. This whitepaper centers on \textbf{Engineering-as-Code (EaC) for Quantum Application-Specific Integrated Circuits (Quantum ASICs)}: a code-first methodology in which protocol topology, gate sets, simulation stacks, and algorithmic control are defined and executed as software, with programmable holographic metasurfaces as the physical target. We outline the EaC pipeline---protocol layer (minimal topology, restricted gate set), cryogenic meta-atoms (rf-SQUIDs, lithium niobate BAW), 28\,nm FD-SOI Cryo-CMOS control, and code-driven simulation (scqubits, SuperScreen, QuTiP) and routing (QUBO/QAOA, DNN inverse design). Applications---on-chip quantum bus, satellite communications (SATCOM), terrestrial point-to-point backhaul, and quantum illumination---are presented as outcomes of this methodology. The document is partner-neutral and cites capabilities by type only.
\end{abstract}

\section{Introduction: Scalability Crisis and the EaC Response}

Scaling solid-state quantum processors to fault tolerance is hindered by a severe \emph{wiring crisis}: driving modern superconducting qubits requires numerous physical coaxial interconnects from room-temperature electronics down to the millikelvin stages of a dilution refrigerator\cite{ref1}. This static wiring dissipates active power and conducts passive heat; a single UT-085-SS-SS coaxial cable can introduce nearly 1\,mW of heat load to the 4\,K stage, overwhelming the micro-watt cooling capacities at the lowest cryogenic stages\cite{ref1}. The physical footprint of these interconnects also restricts multi-qubit gates to static planar topologies, limiting nearest-neighbor interactions and preventing dynamic all-to-all connectivity.

Concurrently, distributing quantum entanglement over global distances faces environmental limits. Optical Quantum Key Distribution (QKD) suffers from Mie scattering because optical wavelengths are comparable to atmospheric aerosols; microwave photons avoid this via Rayleigh scattering but are susceptible to ambient thermal noise\cite{ref1}.

To decouple scaling from these thermodynamic and atmospheric constraints, we adopt an \textbf{Engineering-as-Code (EaC) for Quantum ASICs} approach. The hardware is defined and driven by software: protocol topology and gate sets are specified in code; simulation (Hamiltonian diagonalization, 2D London equation solvers, open quantum systems) and algorithmic control (QUBO/QAOA routing, DNN inverse design) run as code-executed workflows. The physical target is a programmable holographic metasurface acting as a software-defined quantum bus, realized with cryogenic solid-state meta-atoms and ultra-low-power Cryo-CMOS control. Applications---on-chip entanglement and teleportation, weather-resilient SATCOM, terrestrial P2P backhaul, and quantum illumination---follow as outcomes of this EaC methodology.

\section{Engineering-as-Code and the Quantum ASIC Paradigm}

\subsection{EaC as Organizing Principle}

Under EaC, every stage of the Quantum ASIC is code-defined: the protocol layer (topology, gates, protocol mapping), the simulation stack (Hamiltonian and electrodynamic solvers, open-system noise models), and the algorithmic control loop (routing optimization, phase synthesis). The pipeline produces structured outputs (e.g., mapping JSON, phase arrays) that would drive control firmware when physical hardware is deployed. The QASIC Engineering-as-Code repository implements this pipeline; see the repository README and the expanded Markdown whitepaper §10 for layout and commands.

\subsection{Protocol Layer and Minimal Quantum ASIC}

The Quantum ASIC is designed as the minimum hardware required for a target set of protocols---entanglement distribution, teleportation, tamper-evident links, bit commitment---rather than universal quantum computation. The logical architecture is defined in code (e.g., \texttt{asic/topology.py}, \texttt{asic/gate\_set.py}) before any electromagnetic simulation.

\textbf{Topology:} Exactly three logical qubits in a strict linear chain: $0 - 1 - 2$. Only adjacent pairs $(0,1)$ and $(1,2)$ support native two-qubit gates.

\textbf{Gate set:} Single-qubit: Hadamard (H), Pauli-X (X), Pauli-Z (Z), and parameterized $R_x(\theta)$. Two-qubit: Controlled-NOT (CNOT) only, on edges $(0,1)$ or $(1,2)$. All-to-all connectivity, SWAP, T, and Toffoli are excluded; the metasurface routing layer bridges non-adjacent interactions holographically.

\begin{table}[htbp]
\centering
\renewcommand{\arraystretch}{1.4}
\begin{tabularx}{\textwidth}{>{\raggedright\arraybackslash}p{2.8cm}>{\raggedright\arraybackslash}X>{\raggedright\arraybackslash}X}
\toprule
\textbf{Protocol} & \textbf{Qubit roles} & \textbf{ASIC gate sequence} \\
\midrule
Quantum Teleportation & 0: message; 1,2: Bell pair (Alice, Bob) & H(1), CNOT(1,2), CNOT(0,1), H(0) \\
Bit Commitment & 0: Alice; 1: Bob & H(0), CNOT(0,1); optional X(0) to commit-to-1 \\
Tamper-Evident (Thief) & Same as Teleportation & Teleport ops + $R_x(\theta)$ on qubit 2 \\
\bottomrule
\end{tabularx}
\caption{Protocol validations in the software-defined Quantum ASIC layer.}
\end{table}

State-vector simulations (e.g., \texttt{asic/executor.py}) validate these sequences in silico before mapping to physical microwave pulses.

\section{Cryogenic Hardware: Target of the EaC Pipeline}

The EaC stack designs and controls physical components that must operate at 10\,mK. Microwave photon energy $E = h\nu$ is minute; the metasurface control layer must operate well below 100\,mK to preserve coherence and avoid \emph{thermal occlusion}, where $k_B T$ masks the quantum signal. Nematic liquid crystals and phase-change materials freeze or require Joule heating and are incompatible; the system uses purely solid-state, flux-tunable meta-atoms.

\subsection{Magnetic Meta-Atoms: rf-SQUIDs}

Radio-frequency Superconducting Quantum Interference Devices (rf-SQUIDs) provide flux-tunable kinetic inductance: a single superconducting loop interrupted by one Josephson junction. The phase across the junction is set by external magnetic flux, so the effective inductance can be tuned to steer microwave phase. Dense arrays coupled to niobium transmission lines act as a nonlinear, flux-tunable effective medium. Challenges include \emph{dissipative discrete breathers}---localized excitations that trap energy and disrupt phase gradients. Coherence is restored by tuning DC and RF flux to force the array into a synchronized state. \emph{Reflective-type} topologies (hybrid coupler with rf-SQUID reflective loads) maintain 50-$\Omega$ impedance and broadband true-time delay while avoiding the impedance mismatch of direct transmissive phase shifters.

\subsection{Acoustic Meta-Atoms: Lithium Niobate BAW}

Piezoelectric lithium niobate ($\mathrm{LiNbO}_3$) Bulk Acoustic Wave (BAW) resonators mediate microwave-to-phonon coupling. At 4\,K, attenuation follows the Landau-Rumer model and quality factors can reach millions; at 10\,mK, Two-Level System (TLS) defects freeze into coherent states and resonantly absorb phonons, degrading $Q$. Mitigation (phononic bandgaps, DC electric field tuning of TLS frequencies) is validated in code. Integration via split-post microwave cavities aligns cavity field maxima with acoustic mode volume, achieving optomechanical coupling rates far larger than with quartz BAWs.

\subsection{28\,nm FD-SOI Cryo-CMOS Control}

Actuating a 1,000-element metasurface with room-temperature DACs and coaxial cables would thermally overload the cryostat. Heterogeneously integrated 28\,nm Fully Depleted Silicon-On-Insulator (FD-SOI) Cryo-CMOS controllers co-located at 10--100\,mK replace thousands of analog lines with a digital interface (e.g., 4-wire SPI). Charge-Lock Fast-Gate (CLFG) cells trap charge on floating capacitors to generate bias voltages without continuous current. Measured active dissipation is about 18\,nW per cell; 1,000 cells thus dissipate about 18\,$\mu$W---far below the heat load of a single coaxial cable and within dilution refrigerator budgets.

\begin{table}[htbp]
\centering
\renewcommand{\arraystretch}{1.3}
\begin{tabular}{p{0.28\textwidth}p{0.32\textwidth}p{0.32\textwidth}}
\toprule
\textbf{Aspect} & \textbf{Traditional coax routing} & \textbf{28\,nm FD-SOI Cryo-CMOS} \\
\midrule
Interconnects & 1 coax per meta-atom & 4-wire digital SPI \\
Heat load (4\,K/10\,mK) & $\sim$1\,mW per line & 18\,nW per cell \\
Signal generation & Room-temperature DACs & On-chip CLFG capacitors \\
Scalability & Cooling power limit & $>$1,000 meta-atoms \\
\bottomrule
\end{tabular}
\caption{Control architecture comparison.}
\end{table}

\section{Simulation and Computational Stack (EaC Core)}

Under EaC, every simulation stage is code-defined. Python libraries and declarative configurations replace ad-hoc manual analysis.

\subsection{Hamiltonian and Single-Meta-Atom: scqubits}

The rf-SQUID Hamiltonian is
\begin{equation}
H = 4E_C n^2 - E_J \cos(\phi) + \frac{1}{2} E_L(\phi - \phi_{\mathrm{ext}})^2.
\end{equation}
The open-source library \texttt{scqubits} diagonalizes this (and related circuit models) via SciPy backends. Parameter sweeps identify flux ``sweet spots'' where nonlinearity is maximized and sensitivity to $1/f$ flux noise is minimized. The API integrates with QuTiP for composite Hilbert spaces and time evolution.

\subsection{2D Meissner Screening and Array Inductance: SuperScreen}

Scaling to a 1,000-element array makes full 3D FEM intractable for thin-film aspect ratios. \texttt{SuperScreen} solves the 2D London equation by matrix inversion, computing Meissner screening currents and self- and mutual-inductance matrices for aperiodic arrays. The sheet current relates to the vector potential and phase via
\begin{equation}
\mathbf{K}(x,y) = -\frac{1}{\mu_0 \Lambda(x,y)} \left[ \mathbf{A}(x,y) + \frac{\Phi_0}{2\pi} \nabla\theta(x,y) \right].
\end{equation}
This bridges quantum kinetics (e.g., from scqubits) with macroscopic array electrodynamics.

\subsection{Open Quantum Systems and TLS: QuTiP}

At 10\,mK, TLS defects cause dielectric noise and $T_1$/$T_2$ decay. Dynamics are modeled with the Lindblad master equation:
\begin{equation}
\dot{\rho} = -\frac{i}{\hbar}[H_{\mathrm{sys}} + H_{\mathrm{bath}} + H_{\mathrm{int}}, \rho] + \sum_k \left( L_k \rho L_k^\dagger - \frac{1}{2}\{L_k^\dagger L_k, \rho\} \right).
\end{equation}
QuTiP provides \texttt{DecoherenceNoise} and collapse operators to integrate this. Code-based mitigation (e.g., phononic bandgaps, DC tuning of TLS frequencies) can predict coherence improvements of up to two orders of magnitude before fabrication.

\begin{table}[htbp]
\centering
\renewcommand{\arraystretch}{1.35}
\begin{tabularx}{\textwidth}{>{\raggedright\arraybackslash}p{2.2cm}>{\raggedright\arraybackslash}X>{\raggedright\arraybackslash}X}
\toprule
\textbf{Tool} & \textbf{Target} & \textbf{Role in EaC stack} \\
\midrule
scqubits & Energy spectra, coherence & Hamiltonian diagonalization; sweet spots \\
SuperScreen & Meissner currents, $L$, $M$ matrices & 2D London equation; array impedance \\
QuTiP & $T_1$, $T_2$, TLS dynamics & Lindblad master equation; noise mitigation \\
\bottomrule
\end{tabularx}
\caption{Simulation stack for cryogenic metamaterials (EaC core).}
\end{table}

Optional tools (SQDMetal, JosephsonCircuits.jl, BeamProp) support RF scattering, harmonic balance, and acoustic mode propagation where needed.

\section{Algorithmic Control Loop (EaC: Routing and Inverse Design)}

\subsection{QUBO and QAOA Routing}

Mapping logical qubits to physical nodes to minimize gate depth and SWAP count is NP-hard. It is formulated as a Quadratic Unconstrained Binary Optimization (QUBO): binary variables and a cost function that penalizes inefficient wiring. The Quantum Approximate Optimization Algorithm (QAOA) is used as a hybrid quantum-classical solver. The system is initialized in the ground state of a Mixing Hamiltonian $H_M = \sum_i \sigma_i^x$ (all routing configurations in superposition), then alternates evolution under a Problem Hamiltonian $H_P$ (encoding the QUBO cost) and $H_M$ to tunnel out of local minima. Implementations use IBM Qiskit and OpenQAOA; classical optimizers adjust angles $(\beta, \gamma)$. On IBM hardware (e.g., \texttt{ibm\_fez}), the 3-qubit minimal topology has achieved an optimized logical-to-physical mapping with objective penalty 4.0.

\subsection{DNN Inverse Design and Phase Synthesis}

The QAOA output is a low-dimensional routing vector (e.g., logical-to-physical mapping). Translating this into continuous phase commands for 1,000 meta-atoms in real time is infeasible with full-wave FDTD/FEM. The pipeline uses PyTorch-based Deep Neural Network (DNN) surrogates. A forward prediction network (CNN) is trained on FEM-simulated configurations to predict S-parameters and near-field patterns. An inverse design network takes the topology feature vector $\mathbf{T}_{\mathrm{target}} \in \mathbb{R}^d$ and, via hidden layers (e.g., ReLU), produces the phase array:
\begin{equation}
\boldsymbol{\Phi}_{\mathrm{pred}} = \mathrm{Sigmoid}(\mathbf{W}_{\mathrm{out}} \mathbf{H}_n + \mathbf{b}_{\mathrm{out}}) \times 2\pi \;\in\; [0,2\pi]^M.
\end{equation}
Empirically, the generated phases concentrate in a narrow band (e.g., 3.03--3.28\,rad, mean $\approx \pi$). A $\pi$-radian baseline acts as a flat inversion mirror; sub-radian perturbations steer the beam. This constraint keeps Cryo-CMOS actuation within the 18\,nW-per-cell budget and avoids thermal overload.

Routing and inverse design are \emph{code-executed} workflows: Qiskit/OpenQAOA and PyTorch produce mapping JSON and phase arrays (e.g., \texttt{routing\_result.json}, \texttt{phases.npy}) consumed by downstream control or simulation.

\section{Applications of the EaC Quantum ASIC}

The following are \textbf{applications enabled by} the code-defined Quantum ASIC and metasurface control, not separate threads. Each relies on the same EaC pipeline (protocol, simulation, routing, inverse design) and cryogenic hardware.

\subsection{On-Chip Quantum Bus}

The metasurface realizes dynamic all-to-all connectivity on a planar chip without static waveguides. Protocol validation (teleportation, tamper-evident ``Thief'' circuit, bit commitment) runs in simulation and, with hardware, would use the same minimal topology and gate set; the metasurface schedule realizes the required connectivity at each time step.

\subsection{SATCOM: Weather-Resilient Microwave QKD}

Optical QKD is limited by Mie scattering (wavelength $\sim$ droplet size). Microwave photons follow Rayleigh scattering and diffract around aerosols, enabling weather-resilient links. Thermal noise is suppressed by \emph{radiative cooling}: overcoupling the microwave channel to a 10\,mK cold load (or 2.7\,K space) drives the effective thermal photon occupation to about 0.06. Unconditionally secure continuous-variable teleportation can achieve state transfer fidelities of about 58.5\% even at 4\,K ambient.

\subsection{Terrestrial P2P Backhaul}

In terrestrial point-to-point backhaul (e.g., 6G contexts), links face rain attenuation, atmospheric absorption, and Rayleigh/Rician fading. The same Quantum ASIC and EaC methodology apply: radiative sky cooling (e.g., 8--13\,$\mu$m atmospheric window) and 10\,mK overcoupling of transmission lines suppress thermal noise; continuous-variable quantum teleportation (CVQT) with Noiseless Linear Amplification (NLA) and multipartite states (EPR, GHZ, cluster) supports multi-user mesh routing. Environmental impediments are addressed within the same code-defined metasurface and control framework.

\subsection{Quantum Illumination}

The metasurface can generate and steer Two-Mode Squeezed Vacuum (TMSV) states. The signal mode is transmitted into a noisy or cluttered environment; the idler is retained. Although entanglement is broken in transit, \emph{quantum discord} survives. Joint measurement of return signal and idler doubles the quantum Chernoff exponent relative to classical heterodyne receivers, improving detection to the Heisenberg limit. In simulations (e.g., heavy fog, urban clutter), microwave quantum illumination achieves up to 1.5$\times$ range extension and rejects multipath clutter up to 20\,dB stronger than target echoes.

\begin{table}[htbp]
\centering
\renewcommand{\arraystretch}{1.3}
\begin{tabular}{p{0.22\textwidth}p{0.26\textwidth}p{0.28\textwidth}p{0.16\textwidth}}
\toprule
\textbf{Application} & \textbf{Obstacle} & \textbf{Mitigation (EaC/ASIC)} & \textbf{Metric} \\
\midrule
On-chip bus & Static topology & Programmable metasurface routing & Protocol fidelity \\
SATCOM & Mie scattering, thermal noise & Microwave/Rayleigh; radiative cooling & $\bar{n}_{\mathrm{th}} \approx 0.06$; 58.5\% fidelity \\
Terrestrial P2P & Rain, fading, thermal & Radiative cooling; CVQT; NLA; multipartite & Mesh routing; range \\
Quantum Illumination & Clutter, thermal & TMSV; quantum discord; joint measurement & 1.5$\times$ range; 20\,dB clutter rejection \\
\bottomrule
\end{tabular}
\caption{Applications of the EaC Quantum ASIC and mitigation strategies.}
\end{table}

\section{Ecosystem and Cost-Effective Infrastructure (Partner-Neutral)}

Deployment of EaC-designed Quantum ASICs requires access to \emph{types} of capabilities, described here by category only. No specific organization or partnership is implied.

\textbf{Simulation and characterization:} GPU-accelerated HPC for molecular and electromagnetic simulation; low-temperature near-field characterization (e.g., scanning microwave microscopy at millikelvin); integrated photonics packaging and metrology. These capabilities support validation of scqubits/SuperScreen/QuTiP models and refinement of DNN forward/inverse models.

\textbf{Cost-effective compute:} Specialized GPU cloud providers (for PyTorch training and large-scale FEM); spot or preemptible instances (for parameter sweeps, QuTiP trajectories); federated academic grids (for non-proprietary, high-throughput workloads); cloud or local quantum simulators (e.g., \texttt{qiskit-aer}) for QAOA validation before running on physical QPUs.

\section{Roadmap and Implementation}

\textbf{EaC milestones:} (1) Cryogenic metasurface demonstrator---validate tuning at 10--100\,mK with negligible added decoherence; (2) on-chip entanglement via bus---Bell pair and teleportation with metasurface mediation; (3) protocol-layer validation---map teleportation and tamper-evident circuits to the minimal gate set and compare fidelities to simulation; (4) link to SATCOM/terrestrial---over-the-air entanglement or teleportation over short then longer links; (5) scaling---larger apertures, more qubits, hybrid quantum-classical payloads.

\textbf{Implementation:} The QASIC Engineering-as-Code repository provides protocol validation and the full EaC pipeline: QUBO/QAOA routing (simulation or real IBM hardware), inverse design (topology $\to$ phase profile), and pipeline/visualization scripts. Outputs (e.g., \texttt{routing\_result.json}, \texttt{phases.npy}, inverse JSON) are the same data that would drive control firmware when a physical metasurface is available. See the repository README and the expanded Markdown whitepaper §10 for layout, commands, and output formats.

\section{Conclusions}

The scalability crisis in quantum hardware---wiring heat load and static topologies at millikelvin, plus environmental limits on entanglement distribution---motivates a shift to \textbf{Engineering-as-Code for Quantum ASICs}. Under this methodology, protocol topology and gate sets are defined in software; simulation (scqubits, SuperScreen, QuTiP) and algorithmic control (QUBO/QAOA, DNN inverse design) run as code-executed workflows; and the physical target is a programmable holographic metasurface driven by cryogenic meta-atoms and 28\,nm FD-SOI Cryo-CMOS at 18\,nW per cell. The $\pi$-radian phase baseline from the DNN keeps actuation within thermal budgets. Applications---on-chip quantum bus, weather-resilient SATCOM, terrestrial P2P backhaul, and quantum illumination---are outcomes of the same EaC pipeline. Infrastructure is described by capability type only, with no implied partnership. This unified view provides a single narrative and reference for the EaC-based design and control of Quantum ASICs across the documents in the repository.

\appendix
\section{QAOA and DNN Phase Synthesis (Summary)}

\textbf{QAOA:} Initial state $\ket{\psi_0} = \ket{+}^{\otimes N}$ (all configurations in superposition). Evolution $\ket{\psi(\boldsymbol{\gamma},\boldsymbol{\beta})} = \prod_{k=1}^p e^{-i\beta_k H_M} e^{-i\gamma_k H_P} \ket{\psi_0}$; $H_M$ drives transitions out of local minima. Classical optimizer sets $(\beta,\gamma)$ to minimize the QUBO cost.

\textbf{DNN:} Topology vector $\mathbf{T}_{\mathrm{target}}$ $\to$ hidden layers $\mathbf{H}_l = \sigma(\mathbf{W}_l \mathbf{H}_{l-1} + \mathbf{b}_l)$ $\to$ phase output $\boldsymbol{\Phi}_{\mathrm{pred}} \in [0,2\pi]^M$. Phase band $\approx \pi \pm 0.14$\,rad preserves Cryo-CMOS power budget.

\section{Code and Repository}

Protocol layer: \texttt{state/}, \texttt{protocols/}, \texttt{asic/}. Engineering: \texttt{engineering/routing\_qubo\_qaoa.py}, \texttt{engineering/metasurface\_inverse\_net.py}, \texttt{engineering/run\_pipeline.py}, \texttt{engineering/viz\_routing\_phase.py}. Optional: SuperScreen (\texttt{engineering/superscreen\_demo.py}). Dependencies: \texttt{qiskit}, \texttt{qiskit-optimization}, \texttt{torch}; for real hardware, \texttt{qiskit-ibm-runtime} and \texttt{IBM\_QUANTUM\_TOKEN}.

\begin{thebibliography}{99}

\bibitem{ref1} Internal synthesis: Cryogenic Metamaterial Architectures; Computational Materials Science; Engineering as Code Roadmap; Holographic Metasurfaces whitepaper; quantum-terrestrial-backhaul. 2026.

\bibitem{ref2} Scqubits: a Python package for superconducting qubits. Quantum Science and Technology / arXiv.

\bibitem{ref3} SuperScreen: An open-source package for simulating the magnetic response of two-dimensional superconducting devices. OSTI / documentation.

\bibitem{ref4} QuTiP: Quantum Toolbox in Python. \url{https://qutip.org}

\bibitem{ref5} Qubit Routing for QAOA. OpenQAOA documentation.

\bibitem{ref6} Deep neural networks for the evaluation and design of photonic devices. Nature Reviews / inverse design literature.

\bibitem{ref7} Cryogenic CMOS for qubit control and readout. Nature Electronics / Microsoft Research (Gooseberry).

\bibitem{ref8} Dissipative breathers in rf SQUID metamaterials. arXiv / Semantic Scholar.

\bibitem{ref9} Low Temperature Properties of Low-Loss Macroscopic Lithium Niobate Bulk Acoustic Wave Resonators. arXiv.

\bibitem{ref10} Radiative cooling; 8--13 $\mu$m atmospheric window. OSTI / PMC / Stanford.

\bibitem{ref11} Microwave quantum illumination; quantum discord; Chernoff exponent. Preprints / ResearchGate / literature.

\bibitem{ref12} Continuous-variable quantum teleportation; NLA; multipartite states. arXiv / MDPI.

\bibitem{ref13} Rayleigh and Rician fading; terrestrial propagation. MDPI / Wikipedia / channel models.

\end{thebibliography}

\end{document}
